% GENERATED FILE. Edit canonical lessons, then run npm run generate:books.
% canonical-source-hash: fnv1a64:00038ecc2bed78bb
% canonical-lessons: FA-C354-kashf-kardan, FA-C354-ekhtera-kardan, FA-C354-tasis-kardan, FA-C354-katibe, FA-C354-virane

\chapter{Discoveries}
\label{ch:fa-discoveries}
\hlchaptermodality{354}

\begin{chapteropening}
\textbf{By the end of this chapter:} \emph{I can say what someone discovered, invented or founded, and talk about inscriptions and ruins.}

Complete the last lesson of chapter 354: I can say what someone discovered, invented or founded, and talk about inscriptions and ruins.
\end{chapteropening}

\section[\texorpdfstring{\fa{کشف} \fa{کردن} (kashf kardan)}{کشف کردن (kashf kardan)}]{\fa{کشف} \fa{کردن} (kashf kardan) — to discover}

\label{lesson:FA-C354-kashf-kardan}



\begin{quote}
Before the new one: say the Persian for ancestors, then the Persian for heritage.
\end{quote}

\subsection*{You'll want to know: \fa{کشف} \fa{کردن}}
\textbf{\fa{کشف} \fa{کردن}} (\emph{kashf kardan}) — \textquotedblleft{}to discover\textquotedblright{}.

\begin{quote}
\fa{دانشمندان} \fa{یک} \fa{شهر} \fa{قدیمی} \fa{را} \fa{زیر} \fa{زمین} \fa{کشف} \fa{کردند}.
\end{quote}

\begin{quote}
\emph{Dâneshmandân yek shahr-e qadimi râ zir-e zamin kashf kardand.}
\end{quote}

\begin{quote}
Scientists discovered an old city underground.
\end{quote}

Present \emph{kashf mikonad}. \emph{Kashf} alone is \textquotedblleft{}a discovery\textquotedblright{}.

\begin{cousinweb}
\emph{kashf} is borrowed from Arabic \emph{kashf}, \textquotedblleft{}an unveiling, an uncovering\textquotedblright{} (EtymDB). \emph{kardan} \textquotedblleft{}to do\textquotedblright{} is inherited Persian.
\end{cousinweb}

\begin{grammarlens}[title={What you've built}]
One more word for telling, explaining and describing at B1.
\end{grammarlens}

\subsection*{Your turn}
\begin{itemize}
\raggedright
  \item \emph{Say it:} \emph{kashf kardan}
  \item \emph{Say it:} \emph{kashf kardan}, once more
  \item \emph{Say it:} a sentence of your own with \emph{kashf kardan}
  \item \emph{Recall:} say the Persian for ancestors, then the Persian for heritage, then say \emph{kashf kardan} again
\end{itemize}

\begin{tcolorbox}[breakable,colback=teal!4,colframe=teal!35!black,title={Before you move on}]
What does \fa{کشف} \fa{کردن} mean? (\textquotedblleft{}To discover\textquotedblright{}.) Say it once more.
\end{tcolorbox}
\section[\texorpdfstring{\fa{اختراع} \fa{کردن} (ekhterâ' kardan)}{اختراع کردن (ekhterâ' kardan)}]{\fa{اختراع} \fa{کردن} (ekhterâ' kardan) — to invent}

\label{lesson:FA-C354-ekhtera-kardan}



\begin{quote}
Before the new one: say the Persian for heritage, then the Persian for to discover.
\end{quote}

\subsection*{You'll want to know: \fa{اختراع} \fa{کردن}}
\textbf{\fa{اختراع} \fa{کردن}} (\emph{ekhterâ' kardan}) — \textquotedblleft{}to invent\textquotedblright{}.

\begin{quote}
\fa{او} \fa{یک} \fa{قفل} \fa{جدید} \fa{اختراع} \fa{کرد}.
\end{quote}

\begin{quote}
\emph{U yek qofl-e jadid ekhterâ' kard.}
\end{quote}

\begin{quote}
He invented a new kind of lock.
\end{quote}

Present \emph{ekhterâ' mikonad}. \emph{Ekhterâ'} alone is \textquotedblleft{}an invention\textquotedblright{}.

\begin{cousinweb}
\emph{ekhterâ'} is borrowed from Arabic \emph{ikhtirā'}, \textquotedblleft{}invention\textquotedblright{} (EtymDB lists the Arabic word). \emph{kardan} \textquotedblleft{}to do\textquotedblright{} is inherited Persian.
\end{cousinweb}

\begin{grammarlens}[title={What you've built}]
One more word for telling, explaining and describing at B1.
\end{grammarlens}

\subsection*{Your turn}
\begin{itemize}
\raggedright
  \item \emph{Say it:} \emph{ekhterâ' kardan}
  \item \emph{Say it:} \emph{ekhterâ' kardan}, once more
  \item \emph{Say it:} a sentence of your own with \emph{ekhterâ' kardan}
  \item \emph{Recall:} say the Persian for heritage, then the Persian for to discover, then say \emph{ekhterâ' kardan} again
\end{itemize}

\begin{tcolorbox}[breakable,colback=teal!4,colframe=teal!35!black,title={Before you move on}]
What does \fa{اختراع} \fa{کردن} mean? (\textquotedblleft{}To invent\textquotedblright{}.) Say it once more.
\end{tcolorbox}
\section[\texorpdfstring{\fa{تاسیس} \fa{کردن} (ta'sis kardan)}{تاسیس کردن (ta'sis kardan)}]{\fa{تاسیس} \fa{کردن} (ta'sis kardan) — to found, to establish}

\label{lesson:FA-C354-tasis-kardan}



\begin{quote}
Before the new one: say the Persian for to discover, then the Persian for to invent.
\end{quote}

\subsection*{You'll want to know: \fa{تاسیس} \fa{کردن}}
\textbf{\fa{تاسیس} \fa{کردن}} (\emph{ta'sis kardan}) — \textquotedblleft{}to found, to establish\textquotedblright{}.

\begin{quote}
\fa{پدربزرگم} \fa{این} \fa{مدرسه} \fa{را} \fa{تاسیس} \fa{کرد}.
\end{quote}

\begin{quote}
\emph{Pedarbozorgam in madrese râ ta'sis kard.}
\end{quote}

\begin{quote}
My grandfather founded this school.
\end{quote}

Careful print adds a hamza sign. Present \emph{ta'sis mikonad}.

\begin{cousinweb}
\emph{ta'sis} is borrowed from Arabic \emph{ta'sīs}, \textquotedblleft{}a founding, laying a foundation\textquotedblright{} (EtymDB lists the Arabic word). \emph{kardan} \textquotedblleft{}to do\textquotedblright{} is inherited Persian.
\end{cousinweb}

\begin{grammarlens}[title={What you've built}]
One more word for telling, explaining and describing at B1.
\end{grammarlens}

\subsection*{Your turn}
\begin{itemize}
\raggedright
  \item \emph{Say it:} \emph{ta'sis kardan}
  \item \emph{Say it:} \emph{ta'sis kardan}, once more
  \item \emph{Say it:} a sentence of your own with \emph{ta'sis kardan}
  \item \emph{Recall:} say the Persian for to discover, then the Persian for to invent, then say \emph{ta'sis kardan} again
\end{itemize}

\begin{tcolorbox}[breakable,colback=teal!4,colframe=teal!35!black,title={Before you move on}]
What does \fa{تاسیس} \fa{کردن} mean? (\textquotedblleft{}To found, to establish\textquotedblright{}.) Say it once more.
\end{tcolorbox}
\section[\texorpdfstring{\fa{کتیبه} (katibe)}{کتیبه (katibe)}]{\fa{کتیبه} (katibe) — an inscription}

\label{lesson:FA-C354-katibe}



\begin{quote}
Before the new one: say the Persian for to invent, then the Persian for to found.
\end{quote}

\subsection*{You'll want to know: \fa{کتیبه}}
\textbf{\fa{کتیبه}} (\emph{katibe}) — \textquotedblleft{}an inscription\textquotedblright{}.

\begin{quote}
\fa{روی} \fa{سنگ} \fa{یک} \fa{کتیبه} \fa{قدیمی} \fa{بود}.
\end{quote}

\begin{quote}
\emph{Ru-ye sang yek katibe-ye qadimi bud.}
\end{quote}

\begin{quote}
There was an old inscription on the stone.
\end{quote}

Carved on stone or set in tiles on a building.

\begin{cousinweb}
\emph{katibe} is usually traced to Arabic \emph{katība} (EtymDB lists the Arabic word), built on the root \emph{k-t-b} \textquotedblleft{}to write\textquotedblright{}, the root of \emph{ketâb} \textquotedblleft{}book\textquotedblright{}.
\end{cousinweb}

\begin{grammarlens}[title={What you've built}]
One more word for telling, explaining and describing at B1.
\end{grammarlens}

\subsection*{Your turn}
\begin{itemize}
\raggedright
  \item \emph{Say it:} \emph{katibe}
  \item \emph{Say it:} \emph{katibe}, once more
  \item \emph{Say it:} a sentence of your own with \emph{katibe}
  \item \emph{Recall:} say the Persian for to invent, then the Persian for to found, then say \emph{katibe} again
\end{itemize}

\begin{tcolorbox}[breakable,colback=teal!4,colframe=teal!35!black,title={Before you move on}]
What does \fa{کتیبه} mean? (\textquotedblleft{}An inscription\textquotedblright{}.) Say it once more.
\end{tcolorbox}
\section[\texorpdfstring{\fa{ویرانه} (virâne)}{ویرانه (virâne)}]{\fa{ویرانه} (virâne) — a ruin, ruins}

\label{lesson:FA-C354-virane}



\begin{quote}
Before the new one: say the Persian for to found, then the Persian for an inscription.
\end{quote}

\subsection*{You'll want to know: \fa{ویرانه}}
\textbf{\fa{ویرانه}} (\emph{virâne}) — \textquotedblleft{}a ruin, ruins\textquotedblright{}.

\begin{quote}
\fa{از} \fa{قلعه} \fa{فقط} \fa{یک} \fa{ویرانه} \fa{مانده} \fa{است}.
\end{quote}

\begin{quote}
\emph{Az qal'e faqat yek virâne mânde ast.}
\end{quote}

\begin{quote}
Only a ruin is left of the castle.
\end{quote}

From \emph{virân} \textquotedblleft{}ruined\textquotedblright{} + \emph{-e}.

\begin{cousinweb}
\emph{virâne} adds the suffix \emph{-e} to \emph{virân} \textquotedblleft{}desolate, ruined\textquotedblright{}, a native Persian word that EtymDB lists without a parent.
\end{cousinweb}

\begin{grammarlens}[title={What you've built}]
One more word for telling, explaining and describing at B1.
\end{grammarlens}

\subsection*{Your turn}
\begin{itemize}
\raggedright
  \item \emph{Say it:} \emph{virâne}
  \item \emph{Say it:} \emph{virâne}, once more
  \item \emph{Say it:} a sentence of your own with \emph{virâne}
  \item \emph{Recall:} say the Persian for to found, then the Persian for an inscription, then say \emph{virâne} again
\end{itemize}

\begin{tcolorbox}[breakable,colback=teal!4,colframe=teal!35!black,title={Before you move on}]
What does \fa{ویرانه} mean? (\textquotedblleft{}A ruin, ruins\textquotedblright{}.) Say it once more.
\end{tcolorbox}
